\subsection*{The actual unrefined product is closed}

\statement{B.52}{Proved}{The tetrahedral--cyclic product is a finite word-equation locus}
Let $\Gamma=\Gamma_{0,[12]}$, let $P$ be its pure subgroup,
and retain the icosahedral embedding $\iota:A_5\hookrightarrow\Gamma$.
Choose $T\cong A_4$ and $D=\langle d\rangle\cong C_5$ in $A_5$,
and put $H_T=C_P(T)$ and $H_5=C_P(D)$.
Then $H_TH_5$ is closed in the full profinite topology of $P$.
The full centralizer product $C_\Gamma(T)C_\Gamma(D)$ is
closed in the full profinite topology of $\Gamma$.
So is $F_TF_5$ for the explicitly normalized rank-two
subgroup $F_5$ constructed below, with
$F_T=C_\Gamma(T)$ as in B.30.
More precisely, choose generators $t_1,t_2$ of $T$ and any
finite presentation
\[
 A_5=\langle t_1,t_2,d\mid\mathcal R\rangle.
\]
With these letters interpreted as the fixed elements of $\Gamma$,
\[
 H_TH_5
 =\{g\in P:R(t_1,t_2,gdg^{-1})=1
                    \text{ for every }R\in\mathcal R\}.
 \tag{B.52.1}
\]
Consider embeddings $A_5\hookrightarrow\Gamma$ with the
same elementwise permutation action as $\iota$ on the twelve
punctures.  They form exactly four $P$-conjugacy classes.
A residue in $(\mathbb Z/5)^\times$ at one named $D$-fixed
puncture distinguishes these four classes.
These are the exact unrefined products proved for route B.4(c).
It does not prove closure of all its finite-index refinements;
the exact remaining obstruction is given below.

\begin{proof}
Identify the twelve punctures with the transitive $A_5$-set
$A_5/D$, and write $p=D$ for the named puncture and
$\rho:A_5\hookrightarrow S_{12}$ for this permutation action.

\emph{Unlabelled realization and the discrete normalizer.}
Every subgroup of $\Gamma$ isomorphic to $A_5$ is conjugate
to $\iota(A_5)$.  This is the $r=12$, $N=A_5$ case of
\cite[Theorem 4]{DE_Stukow2004}.  Its geometric proof in this
case is especially short.  Nielsen realization, followed by
filling the punctures, realizes the finite group by M\"obius
transformations of the Riemann sphere.  The icosahedral action
has exceptional orbit sizes $12,20,30$, and all other orbits
have size $60$.  Thus the twelve marked points must be its
unique orbit of size twelve.  The standard conjugacy of
icosahedral rotation groups now gives the required conjugacy
of marked sphere actions.  Nielsen realization in the
punctured setting is supplied by
\cite[Section IV, Theorem 5]{DE_Kerckhoff1983}.

The discrete normalizer is
\[
 N_\Gamma(\iota(A_5))=\iota(A_5),\qquad
 C_\Gamma(\iota(A_5))=1.
 \tag{B.52.2}
\]
For completeness, the regular covering comparison used in
B.7 embeds the normalizer modulo the deck group in
the quotient mapping class group.
The three distinguished points have distinct orders
$2,3,\infty$, so all are fixed individually.  That
three-point pure mapping class group is trivial.
The centralizer is then the center of $A_5$, which is trivial.
Only the discrete covering comparison is used here.

\emph{A rotation residue defined by a central extension.}
Let $\Gamma_p\leq\Gamma$ fix the named puncture $p$.
Capping a disk with eleven punctures gives the exact sequence
\[
 1\longrightarrow\langle Z\rangle
 \longrightarrow B_{11}\longrightarrow\Gamma_p
 \longrightarrow1,
 \qquad Z=\Delta_{11}^{\,2}.
 \tag{B.52.3}
\]
Here $Z$ is the positive boundary twist, and $B_{11}$ is
torsion-free; both facts, including the capping identification,
are proved in \cite[Section 2.2]{DE_GonzalezMeneses2011}.
If $f\in\Gamma_p$ has order five, choose a lift $\beta$.
Write $\beta^5=Z^m$ and define
\[
 \kappa_p(f)=m\pmod5.
\]
Changing the lift to $\beta Z^n$ changes $m$ to $m+5n$,
so the residue is well defined.  If $m=5n$, the element
$\beta Z^{-n}$ would be torsion, and hence the identity;
then $f=1$, a contradiction.  Thus $\kappa_p(f)\ne0$.
Centrality of $Z$ gives
\[
 \kappa_p(cfc^{-1})=\kappa_p(f)\quad(c\in\Gamma_p),
 \qquad
 \kappa_p(f^r)=r\kappa_p(f)\quad(5\nmid r).
 \tag{B.52.4}
\]
In particular no choice of a representative finite-order
homeomorphism, and no uniqueness assertion for its local
angle, is needed for this invariant.

\emph{The four labelled embedding classes.}
Let $\varphi:A_5\hookrightarrow\Gamma$ satisfy
$\operatorname{perm}\circ\varphi=\rho$.
Unlabelled conjugacy gives
\[
 \varphi(a)=k\,\iota(\theta(a))\,k^{-1}
 \quad(a\in A_5),
 \qquad k\in\Gamma,\quad\theta\in\operatorname{Aut}(A_5).
\]
Multiply $k$ on the right by a suitable element of
$\iota(A_5)$ so that $k$ fixes $p$; transitivity on the
punctures permits this.  Replace $\theta$ by the corresponding
inner translate.  Since the stabilizer of $p$ in $A_5$ is
exactly $D$, the elementwise permutation condition now implies
$\theta(D)=D$.  Hence $\theta(d)=d^r$ for a unique
$r\in\{1,2,3,4\}$, and (B.52.4) gives
\[
 \kappa_p(\varphi(d))=r\kappa_p(\iota(d)).
 \tag{B.52.5}
\]

Each value of $r$ occurs.  In the five-letter model take
$d=(0\,1\,2\,3\,4)$ and let $s_r$ be the permutation
$i\mapsto ri\pmod5$.  Conjugation by $s_r$ induces
$\theta_r\in\operatorname{Aut}(A_5)$ with
$\theta_r(d)=d^r$.  The permutation
\[
 \sigma_r(gD)=\theta_r^{-1}(g)D
\]
is well defined, fixes $p$, and satisfies
$\sigma_r\rho(\theta_r(a))\sigma_r^{-1}=\rho(a)$.
Choose a mapping class $k_r\in\Gamma_p$ with permutation
$\sigma_r$; the puncture-permutation map is surjective.
Then $\varphi_r=k_r\iota\theta_r k_r^{-1}$ has the required
elementwise permutation action and residue
$r\kappa_p(\iota(d))$.
For $r=1$ take $k_1=1$, so $\varphi_1=\iota$.

Conversely, every $\varphi$ with the value $r$ in (B.52.5)
is $P$-conjugate to $\varphi_r$.
Indeed $\operatorname{Aut}(A_5)=S_5$, and the centralizer
of the five-cycle $d$ in $S_5$ is $D$.
Consequently $\theta\theta_r^{-1}$ is conjugation by
some $b\in D$.  Absorb $\iota(b)$ into $k$, which still
fixes $p$, so that now $\theta=\theta_r$.
The two permutations $\operatorname{perm}(k)$ and $\sigma_r$
intertwine the same transitive $A_5$ actions and fix $p$.
Their quotient is therefore the identity: a permutation
commuting with a transitive action is determined by its
value at one point.  Thus $k k_r^{-1}\in P$, giving the
asserted conjugacy.  Its conjugator is unique by (B.52.2).
The four residues are distinct, so this proves the exact
classification, and in particular
\[
 \begin{split}
 &\operatorname{perm}\circ\varphi=\rho,\quad
   \kappa_p(\varphi(d))=\kappa_p(\iota(d))\\
 &\hspace{25mm}\Longrightarrow
   \varphi=h\iota h^{-1}\text{ for a unique }h\in P.
 \end{split}
 \tag{B.52.6}
\]

The independent count is
\[
 \frac{|\operatorname{Aut}(A_5)|\,
        |C_{S_{12}}(\rho(A_5))|}
      {|N_\Gamma(\iota(A_5))|}
 =\frac{120\cdot2}{60}=4.
\]
Here the permutation centralizer is
$N_{A_5}(D)/D\cong C_2$.
For each automorphism there are two compatible puncture
permutations, and changing the description by the discrete
normalizer acts freely on these $240$ pairs, with orbits
of size sixty.  The exact program
\texttt{check\_marked\_a5.py} verifies this action and its
four orbits directly, as well as their four rotation
multipliers.  This finite computation supplements the
classification proof; it does not replace the mapping-class
or braid arguments.

\emph{The product as a word-equation locus.}
Suppose $g\in P$ satisfies all the equations in (B.52.1).
The presentation gives a homomorphism
$\varphi:A_5\to\Gamma$ sending
\[
 t_i\longmapsto t_i\quad(i=1,2),\qquad
 d\longmapsto gdg^{-1}.
\]
It is nontrivial and hence injective, because $A_5$ is simple.
Purity of $g$ gives
$\operatorname{perm}\circ\varphi=\rho$ on the three
generators and therefore on the whole group.
Also $g$ fixes $p$, so (B.52.4) gives the equality of
residues in (B.52.6).
Let $h\in P$ be its unique conjugator.
Since $\varphi$ fixes $T$ pointwise, $h\in H_T$;
since $hdh^{-1}=gdg^{-1}$, one has
$h^{-1}g\in H_5$.  Thus $g\in H_TH_5$.
Conversely, if $g=hj$ with $h\in H_T$ and $j\in H_5$,
the generator substitutions in (B.52.1) are precisely
conjugation of $\iota$ by $h$.  All the relators vanish.
This proves (B.52.1).

The certificate makes the presentation fully explicit.
Use $t_1=(0\,1\,2)$, $t_2=(0\,1)(2\,3)$ and the
above $d$ in $A_5$.  Breadth-first search in their Cayley
graph supplies a word $w_a$ for each of the sixty elements,
with $w_1$ empty.  The $180$ relators
\[
 w_a t\,w_{at}^{-1},
 \qquad a\in A_5,\quad t\in\{t_1,t_2,d\},
\]
present $A_5$: they reduce every positive word to its
canonical $w_a$, and make each generator have its finite
permutation order, so the same holds for arbitrary words.
There are at most sixty elements and a surjection to $A_5$.
All words, relators and finite permutation data are recorded
in \texttt{marked\_a5\_certificate.json}.
The relator list is also in \texttt{marked\_a5\_relators.txt}.
Actual standard output is in \texttt{check\_marked\_a5.out}.

Finally, if $g\notin H_TH_5$, some evaluated relator is
a nonidentity element of $\Gamma$.
Residual finiteness of $\Gamma$ gives a finite quotient
detecting that element.  Its restriction to $P$ separates
$g$ from every solution of the equations, and hence from
$H_TH_5$.  This proves closure in the full profinite
topology.  The argument does not use a profinite
centralizer comparison for $A_5$.

\emph{Full centralizers and a specified rank-two cyclic factor.}
Put $C_T=C_\Gamma(T)$ and $C_D=C_\Gamma(D)$.
Their pure kernels $H_T,H_5$ have finite index.
Choose left coset representatives $a_i$ and right coset
representatives $b_j$, so that
\[
 C_T=\bigcup_i a_iH_T,\qquad C_D=\bigcup_j H_5b_j.
\]
Then
\[
 C_TC_D=\bigcup_{i,j}a_i(H_TH_5)b_j.
\]
The full profinite topology on $P$ is induced from
$\Gamma$, because $P$ has finite index, and $P$ is closed
in that topology.  Thus $H_TH_5$ is closed in $\Gamma$.
Translations and finite unions preserve closure, proving
the assertion for the full centralizers.

To remove the finite-factor ambiguity in the notation $F_5$,
make the following explicit choice.  The $D$ action has
two fixed punctures $p,q$ and two free orbits $A,B$ of size
five.  Let $C_D^{\mathrm{ord}}$ be the subgroup of $C_D$
fixing $p,q$ individually and preserving $A,B$ separately.
Choose a named vertex $v\in A$ and put
\[
 F_5^v=\{c\in C_D^{\mathrm{ord}}:c(v)=v\}.
\]
The regular covering comparison gives
\[
 1\longrightarrow D\longrightarrow C_D^{\mathrm{ord}}
 \longrightarrow P\Gamma_{0,4}\longrightarrow1.
\]
Here the quotient has the four labelled punctures
$\bar p,\bar q,\bar A,\bar B$.
Its cyclic-cover monodromy sends their peripheral
generators to $d,d^{-1},1,1$ respectively.
Every pure mapping class preserves this homomorphism,
because it preserves each peripheral conjugacy class
and $D$ is abelian.  Thus every pure base class lifts,
and its lifts commute with $D$; this proves surjectivity
in the displayed sequence.
Every lift has a unique deck translate fixing $v$.
Equivalently, if $c(v)=d^{e(c)}v$, then
$e:C_D^{\mathrm{ord}}\to\mathbb Z/5$ is a homomorphism
whose restriction to $D$ is an isomorphism.
Its kernel is $F_5^v$, and
\[
 C_D^{\mathrm{ord}}=D\times F_5^v,\qquad
 F_5^v\cong P\Gamma_{0,4}\cong F_2.
\]
This is a chosen geometric complement, not an unspecified
quotient by finite factors.
It contains $H_5$.  In fact $[F_5^v:H_5]=5$:
commutation with $D$ and the fixed vertex $v$ fix all
of $A$; the only remaining puncture permutation is a
cyclic shift of $B$.  All five shifts occur in the
family
\[
 \{0,\infty\}\cup\mu_5\cup
       \{z:z^5=\lambda\},\qquad
 \lambda\in\mathbb P^1\setminus\{0,1,\infty\},
\]
which keeps the first pentagon fixed.  A loop of
$\lambda$ around zero produces a generator of the
shift on the second pentagon.

With $F_5=F_5^v$, B.30 gives
$[F_T:H_T]=12$.
Consequently $F_TF_5$ is a union of sixty translates
of $H_TH_5$ and is closed in $\Gamma$.
For another chosen complement $F_5'$ which does not
contain $H_5$, this finite-union argument does not prove
closure of $F_TF_5'$: its pure kernel is a finite-index
refinement of $H_5$, and that product requires the
additional refinement comparison discussed next.
\end{proof}

Here is the precise remaining condition on a finite-index
refinement.  Write the unique factorization of an element
$g\in H_TH_5$ as $g=h(g)j(g)$.
Uniqueness follows from $H_T\cap H_5=1$.
For every $L\leq_f P$, B.28 and B.49 give
\[
 \begin{split}
 g\in P\cap\overline{(H_T\cap L)H_5}
 \quad\Longleftrightarrow\quad
 &g\in H_TH_5\ \text{ and }\\[-1mm]
 &h(g)\mathcal K\cap\overline L\ne\varnothing,
 \end{split}
 \tag{B.52.7}
\]
where closures are in $\widehat P$ and
$\mathcal K=C_{\widehat P}(A_5)$.
Indeed any profinite conjugator for the same embedding as
$h(g)$ is exactly $h(g)k$ with $k\in\mathcal K$;
it belongs to the required refined first factor precisely
when it belongs to $\overline L$.
In contrast, actual membership in $(H_T\cap L)H_5$ is
the condition $h(g)\in L$.
For $U\trianglelefteq_f P$ and the quotient
$q_U:\widehat P\to P/U$, (B.52.7) becomes
\[
 g\in P\cap\overline{(H_T\cap U)H_5}
 \quad\Longleftrightarrow\quad
 g\in H_TH_5\text{ and }q_U(h(g))\in q_U(\mathcal K).
 \tag{B.52.8}
\]
The four $P$-conjugacy classes above therefore do not settle
the conjugacy problem under each smaller subgroup $L$.
Together, B.46 and B.52 reduce $\mathcal K=1$ exactly to
closure of all the refined products $(H_T\cap L)H_5$.
B.58 proves these refinements when $L$ contains $[P,P]$.
The assertion for arbitrary finite-index $L$ remains Open.
